% Part: incompleteness
% Chapter: incompleteness-provability
% Section: provability-conditions

\documentclass[../../../include/open-logic-section]{subfiles}

\begin{document}

\olfileid{inc}{inp}{prc}

\olsection{$\Th{PA}$కు \usetoken{S}{derivability} షరతులు}

పియానో అంకగణితం, లేదా $\Th{PA}$, అన్ని \tetoken{సూత్రాలకూ}{!!{formula}s} ఆగమన స్వీకృతాలను
చేర్చి $\Th{Q}$ను విస్తరించే సిద్ధాంతం. మరో విధంగా, ప్రతి
\tetoken{సూత్రం}{!!{formula}}~$!A$కు $\Th{Q}$లో కింది రూపంలోని స్వీకృతాలను చేరుస్తాం:
\[
(!A(0) \land \lforall[x][(!A(x) \lif !A(x'))]) \lif \lforall[x][!A(x)]
\]
ఇది నిజానికి ఒక \emph{పథకం}, అంటే అనంతంగా అనేక స్వీకృతాలు, అని గమనించండి
(నిజానికి $\Th{PA}$ పరిమితంగా స్వీకృతీకరించదగినది \emph{కాదు}). కానీ ఒక
సంకేతాల వరుస ఆగమన స్వీకృతానికి నిదర్శనమో కాదో ప్రభావవంతంగా నిర్ణయించవచ్చు
కాబట్టి, $\Th{PA}$ స్వీకృతాల సమితి గణనీయమైనది. $\Th{PA}$ అనేది~$\Th{Q}$
కంటే ఎంతో దృఢమైన సిద్ధాంతం. ఉదాహరణకు, సాధారణ విధంగా ఆగమనాన్ని వాడి కూడిక,
గుణకారం పరస్పర వినిమేయాలని సులభంగా నిరూపించవచ్చు. నిజానికి చాలా పరిసమాప్త
సంఖ్యా-సిద్ధాంత, సంయోజనశాస్త్ర వాదనలను~$\Th{PA}$లో నిర్వహించవచ్చు.

$\Th{PA}$ గణనీయంగా స్వీకృతీకరించబడింది కాబట్టి, \tetoken{వ్యుత్పాద్యత}{!!{derivability}}
విధేయం $\Prf[\Th{PA}](x,y)$ గణనీయమైనది; అందువల్ల అది~$\Th{Q}$లో (కాబట్టి
$\Th{PA}$లో కూడా) ప్రాతినిధ్యం చేయబడుతుంది. మునుపటిలాగే, ఆ సంబంధానికి
ప్రాతినిధ్యం వహించే సూత్రాన్ని $\OPrf[\Th{PA}](x,y)$తో సూచిస్తాం.
$\OProv[\Th{PA}](y)$ను
$\lexists[x][\OPrf[\Th{PA}](x,y)]$ అనే సూత్రంగా తీసుకుందాం;
ఇది సహజార్థంలో ``$y$ అనేది $\Th{PA}$ స్వీకృతాల నుంచి
\tetoken{వ్యుత్పాదించదగినది}{!!{derivable}}'' అని చెబుతుంది.\sourcecorrection{OLTEINP-001}{ఆంగ్ల మూలం ఈ నిర్వచనంలో సంబంధ సంకేతం $\Prf$ను వాడింది; ముందరి వాక్యం నిర్వచించిన ప్రాతినిధ్య సూత్రం $\OPrf$ కావాలి. ఆ ఉద్దేశిత సూత్ర సంకేతాన్ని ఇక్కడ నిలిపాం.}
$\Th{Q}$ స్వీకృతాల కంటే కొద్దిగా ఎక్కువ అవసరమవడానికి కారణం, మనం వాడే
సిద్ధాంతం ఈ \tetoken{వ్యుత్పాద్యత}{!!{derivability}} విధేయం గురించిన కొన్ని ప్రాథమిక వాస్తవాలను
\tetoken{నిగమించగలిగేంత}{!!{derive}} బలంగా ఉందని తెలుసుకోవాలి. మనకు కావలసిన వాస్తవాలు ఇవి:
\begin{enumerate}
\item[P1.] $\Th{PA} \Proves !A$ అయితే, $\Th{PA} \Proves
  \OProv[\Th{PA}](\gn{!A})$.
\item[P2.] అన్ని \tetoken{సూత్రాలు}{!!{formula}s} $!A$, $!B$లకు,
  \[
  \Th{PA} \Proves \OProv[\Th{PA}](\gn{!A \lif !B}) \lif
  (\OProv[\Th{PA}](\gn{!A}) \lif \OProv[\Th{PA}](\gn{!B})).
  \]
\item[P3.] ప్రతి \tetoken{సూత్రం}{!!{formula}}~$!A$కు,
  \[
  \Th{PA} \Proves \OProv[\Th{PA}](\gn{!A})
  \lif \OProv[\Th{PA}](\gn{\OProv[\Th{PA}](\gn{!A})}).
  \]
\end{enumerate}
ఈ మూడు లక్షణాలు నిజమని నిర్ధారించే ఏకైక మార్గం
$\OProv[\Th{PA}](y)$ అనే \tetoken{సూత్రాన్ని}{!!{formula}} జాగ్రత్తగా వివరించి, సంబంధిత అధికారిక
\tetoken{వ్యుత్పత్తులను}{!!{derivation}s} వర్ణించేందుకు $\Th{PA}$ స్వీకృతాలను వాడటం. షరతులు (1),
(2) సులభం; నిజంగా శ్రమ కోరేది షరతు~(3). (దానికి ఎలాంటి పని అవసరమో
ఆలోచించండి \dots) వివరాలను పూర్తిగా ఇవ్వడం విసుగ్గా, ఆసక్తిలేకుండా ఉంటుంది;
కాబట్టి $\Th{PA}$కు పైన పేర్కొన్న మూడు లక్షణాలు ఉన్నాయని ఇక్కడ మీరు
నమ్మాలని కోరుతాం. $\OProv[\Th{PA}](y)$కు సమంజసమైన ఎంపిక కింది దానినీ
తృప్తిపరుస్తుంది:
\begin{enumerate}
\item[P4.] $\Th{PA} \Proves \OProv[\Th{PA}](\gn{!A})$ అయితే,
  $\Th{PA} \Proves !A$.
\end{enumerate}
కానీ ఈ వాస్తవం మనకు అవసరం లేదు.

\begin{digress}
యాదృచ్ఛికంగా, ఇప్పుడు మనం చూపుతున్న అలసత్వాన్నే గ్యోడెల్ కూడా చూపించాడు.
1931 పత్రం చివర రెండవ అసంపూర్ణతా సిద్ధాంతపు నిరూపణను స్థూలంగా ఇచ్చి,
వివరాలను తర్వాతి పత్రంలో ఇస్తానని మాటిచ్చాడు. కానీ ఎప్పుడూ ఇవ్వలేదు; ఆ వాదనను
అర్థం చేసుకున్న ప్రతి ఒక్కరూ దానిని పూర్తిచేయవచ్చని నమ్మారు కాబట్టి, ఆయన
వివరాలను పూరించాల్సిన అవసరం రాలేదు.
\end{digress}

\end{document}
